% BE-03: Technology Selection — compile with XeLaTeX (Overleaf: Menu → Compiler → XeLaTeX)
\documentclass[11pt,a4paper]{article}

% ---------- language and fonts ----------
\usepackage{fontspec}
\usepackage{polyglossia}
\setdefaultlanguage{russian}
\setotherlanguage{english}

% pick the first available serif font with Cyrillic
\IfFontExistsTF{Times New Roman}{\setmainfont{Times New Roman}}{%
  \IfFontExistsTF{Liberation Serif}{\setmainfont{Liberation Serif}}{\setmainfont{DejaVu Serif}}}
\IfFontExistsTF{Arial}{\setsansfont{Arial}}{%
  \IfFontExistsTF{Liberation Sans}{\setsansfont{Liberation Sans}}{\setsansfont{DejaVu Sans}}}
\IfFontExistsTF{Consolas}{\setmonofont[Scale=0.9]{Consolas}}{\setmonofont[Scale=0.85]{DejaVu Sans Mono}}

% ---------- layout ----------
\usepackage[margin=2.2cm]{geometry}
\usepackage{microtype}
\usepackage{array}
\usepackage{booktabs}
\usepackage{xltabular}
\usepackage{enumitem}
\usepackage[table]{xcolor}
\usepackage{caption}
\usepackage{fancyvrb}
\usepackage{float}
\usepackage{tikz}
\usetikzlibrary{positioning,arrows.meta,fit,backgrounds}
\usepackage{xurl}
\usepackage[colorlinks=true,linkcolor=blue!50!black,urlcolor=blue!50!black]{hyperref}

\setlist{itemsep=2pt,topsep=4pt}
\setlength{\parskip}{4pt}
\setlength{\parindent}{0pt}
\setlength{\emergencystretch}{3em} % avoids overfull lines if hyphenation patterns are missing
\renewcommand{\arraystretch}{1.25}
\newcolumntype{L}{>{\raggedright\arraybackslash}X}
\newcolumntype{P}[1]{>{\raggedright\arraybackslash}p{#1}}
\newcommand{\tblstart}{\small\rowcolors{2}{gray!7}{white}}
\newcommand{\code}[1]{\texttt{#1}}

\title{\textbf{BE-03: Technology Selection}\\[4pt]
\large Intelligent Kanban System с AI-ассистентом\\ для управления задачами голосом и текстом}
\author{Команда проекта}
\date{Версия 3 · 20 октября 2026}

\begin{document}
\maketitle

\begin{center}
\begin{minipage}{0.9\textwidth}
\small
\textbf{Изменения относительно версии 2:} добавлены frontend, распознавание намерений для 8 AI-действий, права доступа, уведомления, хранение и защита голосовых сообщений, аудит. Решение «локально или облако» пересмотрено с точки зрения реального продукта: скорость, масштабируемость, качество казахского языка.
\end{minipage}
\end{center}

\tableofcontents
\newpage

% =====================================================
\section{Цель}

Выбрать и обосновать технологии для \textbf{коммерческого Kanban-продукта} для сотрудников компаний, руководителей проектов, преподавателей и студентов. AI-ассистент встроен в рабочий процесс: пользователь говорит или пишет команду, система её понимает, показывает предлагаемое действие и после подтверждения выполняет его на доске.

% =====================================================
\section{Требования, определяющие выбор}

{\tblstart
\begin{xltabular}{\textwidth}{P{0.9cm} P{6.2cm} L}
\toprule
\textbf{ID} & \textbf{Требование} & \textbf{Влияние на стек} \\
\midrule
\endhead
R1 & Языки kk, ru, en и смешанная речь (kk + ru + английские термины) & Качество STT и LLM на казахском — главный критерий \\
R2 & 8 AI-действий: Create Task, Update Task, Change Status, Assign User, Change Deadline, Change Priority, Search Tasks, Project Summary & LLM с надёжным structured output / tool calling \\
R3 & AI не меняет данные без подтверждения и при неоднозначности; не выдумывает исполнителей, сроки и задачи & Детерминированная проверка на backend, AI только предлагает \\
R4 & Полноценный продукт: auth, workspaces, проекты, участники, задачи, комментарии, история, уведомления & Реляционная БД, RBAC, real-time обновления \\
R5 & Ответ на голосовую команду за несколько секунд & Быстрый STT, асинхронная обработка \\
R6 & Защита пользовательских данных и голосовых & Шифрование, ограниченный срок хранения, контроль доступа \\
R7 & Масштабируемость и развитие продукта & Stateless API, отдельный AI-воркер, замена моделей конфигурацией \\
R8 & Надёжность: ошибки AI, недоступность API, повторные запросы & Fallback-провайдеры, идемпотентность, очереди \\
R9 & Цифровой след разработки (GitLab, YouTrack) & GitLab CI, тесты, история \\
\bottomrule
\end{xltabular}}

% =====================================================
\section{Архитектура}

\begin{figure}[H]
\centering
\begin{tikzpicture}[
  scale=0.92, transform shape,
  font=\small,
  box/.style={draw, rounded corners=3pt, fill=white, align=center, minimum height=1cm, text width=3.5cm},
  db/.style={draw, rounded corners=9pt, fill=gray!12, align=center, minimum height=1cm, text width=2.5cm},
  grp/.style={draw=gray!60, dashed, rounded corners=6pt, inner sep=8pt},
  lbl/.style={font=\footnotesize\bfseries, text=gray!70!black},
  arr/.style={-{Stealth[length=2mm]}, thick},
  tag/.style={font=\footnotesize, fill=white, inner sep=1.5pt}
]
% frontend
\node[box] (ui)   at (0,0)    {Kanban UI (React)};
\node[box] (mic)  at (0,-1.7) {Запись голоса\\MediaRecorder};
\node[box] (conf) at (0,-3.4) {Окно подтверждения\\AI-действия};
% backend
\node[box] (api)   at (5.7,0)    {REST API + WebSocket};
\node[box] (cmd)   at (5.7,-1.7) {AI Command Service\\проверка и резолвинг};
\node[box] (dom)   at (5.7,-3.4) {Доменные сервисы\\Auth/RBAC, задачи};
\node[box] (notif) at (5.7,-5.1) {Уведомления};
% worker
\node[box] (stt) at (11.4,0)    {Speech-to-Text\\Whisper kk/ru/en};
\node[box] (llm) at (11.4,-1.7) {LLM: intent\\+ entities};
% storage
\node[db] (s3)    at (0,-7)    {MinIO / S3\\голосовые};
\node[db] (pg)    at (8.55,-7)  {PostgreSQL};
\node[db] (redis) at (11.4,-4.6) {Redis\\очереди, pub/sub};

\begin{scope}[on background layer]
\node[grp, fill=blue!4,   fit=(ui)(mic)(conf),        label={[lbl]above:Frontend}] {};
\node[grp, fill=green!5,  fit=(api)(cmd)(dom)(notif), label={[lbl]above:Backend (FastAPI)}] {};
\node[grp, fill=orange!7, fit=(stt)(llm),             label={[lbl]above:AI-воркер}] {};
\end{scope}

% arrows
\draw[arr, <->] (ui) -- node[tag, above=2pt, align=center]{HTTPS\\WSS} (api);
\draw[arr] (api) -- (cmd);
\draw[arr] (conf.east) -- node[tag, sloped, below=1pt, pos=0.4]{подтверждение} (cmd.west);
\draw[arr, <->] (cmd) -- node[tag, above=2pt, align=center]{команда\\план} (llm);
\draw[arr] (stt) -- (llm);
\draw[arr, <->] (llm) -- node[tag, right=2pt]{очередь} (redis);
\draw[arr] (cmd) -- node[tag, right=2pt]{после confirm} (dom);
\draw[arr] (dom) -- (notif);
\draw[arr] (dom.east) -| (pg.north);
\draw[arr] (notif.west) -| node[tag, pos=0.25, below=1pt]{события} (conf.south);
\draw[arr] (mic.west) -- ++(-0.5,0) |- node[tag, pos=0.25, left=1pt]{аудио} (s3.west);
\end{tikzpicture}
\caption{Архитектура компонентов системы. AI-воркер читает аудио из MinIO и возвращает план действий; данные меняет только backend после подтверждения.}
\end{figure}

\subsection{AI-конвейер}

\begin{center}
\small
Voice Input $\rightarrow$ Speech-to-Text $\rightarrow$ Intent Recognition $\rightarrow$ Entity Extraction $\rightarrow$\\
Validation \& Resolution $\rightarrow$ User Confirmation $\rightarrow$ Backend Action $\rightarrow$ Kanban Update
\end{center}

{\tblstart
\begin{xltabular}{\textwidth}{P{3.4cm} P{2.6cm} L}
\toprule
\textbf{Шаг} & \textbf{Где} & \textbf{Технология} \\
\midrule
\endhead
Voice Input & Браузер & MediaRecorder API (WebM/Opus), загрузка в MinIO через backend \\
Speech-to-Text & AI-воркер & Whisper large-v3-turbo + Whisper, дообученный на казахском (раздел~\ref{sec:stt}) \\
Intent Recognition + Entity Extraction & AI-воркер & LLM с tool calling: один вызов возвращает список действий с параметрами (раздел~\ref{sec:llm}) \\
Validation \& Resolution & Backend, без LLM & Pydantic-схемы, поиск исполнителя среди участников проекта, поиск задач (pg\_trgm), проверка дат и прав \\
User Confirmation & Frontend & Карточки-превью с подсвеченными пропущенными и неоднозначными полями; редактирование транскрипта и полей \\
Backend Action & Backend & Те же доменные сервисы, что и для ручных действий в UI (единая проверка прав и аудит) \\
Kanban Update & Frontend & WebSocket-событие, доска обновляется у всех участников \\
\bottomrule
\end{xltabular}}

\textbf{Ключевой принцип.} LLM только \emph{понимает команду и предлагает} действие. Всё, что меняет данные, выполняет обычный backend-код после подтверждения пользователя, с теми же проверками прав, что и для ручных действий. Поэтому AI не может изменить чужую задачу или создать задачу без подтверждения.

% =====================================================
\section{Speech-to-Text}\label{sec:stt}

\subsection{Сравнение}

{\tblstart
\begin{xltabular}{\textwidth}{P{2.9cm} L L L L}
\toprule
\textbf{Критерий} & \textbf{Whisper large-v3-turbo (self-hosted)} & \textbf{Whisper, дообученный на KSC2 (self-hosted)} & \textbf{ElevenLabs Scribe (облако)} & \textbf{OpenAI gpt-4o-transcribe (облако)} \\
\midrule
\endhead
Казахский & Слабо: путает с турецким, в исследовании WER $\approx$77\% & \textbf{Хорошо: WER $\approx$15--18\%} на тесте KSC2 & Хорошо (заявлено WER 3,1\% на FLEURS) & Средне--хорошо \\
Русский / английский & \textbf{Отлично} & Хуже базового & Отлично & Отлично \\
Смешанная речь & Автоопределение & Частично & Автоопределение & Автоопределение \\
Скорость & $\approx$1--2 с на 10 с аудио (GPU) & Аналогично & 1--3 с + сеть & 1--3 с + сеть \\
Стоимость & Только железо & Только железо & $\approx$\$0,22 / час & \$0,006 / мин \\
Приватность & \textbf{Голос не покидает сервер} & \textbf{Голос не покидает сервер} & Данные у провайдера & Данные у провайдера \\
Масштабирование & GPU-воркеры & GPU-воркеры & Автоматически & Автоматически \\
\bottomrule
\end{xltabular}}

\subsection{Решение: self-hosted STT с маршрутизацией по языку и облачный fallback}

\begin{enumerate}
  \item \textbf{Автоматическое определение языка.} Whisper large-v3-turbo (\code{faster-whisper}) вычисляет вероятности языков; они ограничиваются тремя языками (kk, ru, en) и нормируются.
  \item \textbf{Маршрутизация без подсказки.} Казахская модель используется, если казахский — самый вероятный язык \emph{или} $p(\mathrm{kk}) \ge 0{,}30$. Порог нужен для смешанной речи: казахская фраза с русскими и английскими терминами часто определяется как русская, но с заметной долей казахского. В остальных случаях русскую и английскую речь распознаёт Whisper large-v3-turbo с определённым языком.
  \item \textbf{Язык, выбранный пользователем.} Если пользователь выбрал язык в интерфейсе, определение не выполняется: kk — казахская модель, ru/en — Whisper large-v3-turbo. Для смешанной речи интерфейс предлагает казахский.
  \item Пользователь может \textbf{исправить транскрипт} перед обработкой.
  \item \textbf{Fallback:} если локальный STT недоступен или перегружен, используется облачный провайдер (включается в настройках workspace).
\end{enumerate}

Порог $0{,}30$ — начальное значение; он задаётся параметром и подбирается по бенчмарку. Бенчмарк (\code{benchmarks/stt\_benchmark.py}) оценивает многоязычные модели в двух режимах отдельно: \textbf{auto} (автоматическое определение языка — основной сценарий продукта) и \textbf{hint} (язык выбран пользователем). Для режима auto дополнительно считается точность определения языка и то, куда маршрутизируются смешанные записи. Финальная казахская модель выбирается по результатам бенчмарка.

\textbf{Обоснование.} Голос — самые чувствительные данные системы. Self-hosted STT решает вопрос приватности и стоимости при росте числа пользователей, а дообученная модель закрывает слабость Whisper на казахском. Облачный fallback повышает надёжность (R8).

% =====================================================
\section{LLM для распознавания команд}\label{sec:llm}

\subsection{Что должна уметь модель}

Из одной фразы, часто с ошибками распознавания и смешанными языками, получить \textbf{план действий} строго по схеме:

\begin{Verbatim}[fontsize=\small,frame=single,framesep=6pt,rulecolor=\color{gray!50}]
{
  "actions": [
    {
      "type": "create_task",
      "params": {
        "title": "Подготовить презентацию проекта",
        "assignee_mention": "Данияр",
        "deadline": "2026-10-23",
        "deadline_phrase": "до пятницы",
        "priority": "high"
      },
      "missing": [],
      "confidence": 0.92
    }
  ],
  "clarification_needed": null
}
\end{Verbatim}

Требования: tool calling или JSON Schema, понимание kk/ru/en, разбор нескольких поручений в одной фразе, казахские падежные формы имён («Даниярға» $\rightarrow$ «Данияр»), относительные даты.

\subsection{Сравнение}

{\tblstart
\begin{xltabular}{\textwidth}{P{3.2cm} L L}
\toprule
\textbf{Критерий} & \textbf{Локальная LLM 7--8B (Qwen 3 / ISSAI KazLLM, Ollama)} & \textbf{Облачная LLM Flash/mini-класса (Gemini Flash, OpenAI mini)} \\
\midrule
\endhead
Казахский & Qwen — средне; KazLLM — хорошо, но слабее в tool calling & Хорошо \\
Tool calling, несколько действий в одной фразе & Нестабильно на 7--8B, особенно с kk & \textbf{Надёжно} \\
Скорость & 3--15 с без мощной GPU & \textbf{1--3 с} \\
Стоимость & Только железо & Оценка: $\approx$1,5 тыс. входных + 300 выходных токенов на команду $\approx$ \textbf{\$0,005 за команду} \\
Масштабирование & Нужны дополнительные GPU & Автоматически \\
Приватность & Полная & В LLM уходит только текст команды и имена участников проекта; без голоса и содержимого задач \\
\bottomrule
\end{xltabular}}

\subsection{Решение: гибрид через единый интерфейс \code{CommandParser}}

\begin{itemize}
  \item \textbf{Основной режим — облачная LLM Flash/mini-класса} с tool calling. Для реального продукта решают качество разбора сложных команд на казахском, скорость ответа (пользователь ждёт превью карточки) и масштабирование без закупки GPU.
  \item \textbf{Приватный (on-premise) режим — локальная LLM через Ollama} (Qwen 3 или KazLLM 8B). Включается конфигурацией для организаций, которым нельзя отправлять текст во внешние сервисы, и служит fallback при недоступности облачного API.
  \item Оба провайдера реализуют один интерфейс и одну JSON-схему. Конкретная модель выбирается по тесту на наших 30+ сценариях: точность определения намерения, точность полей, доля выдуманных значений, время ответа.
\end{itemize}

\textbf{Почему изменено решение версии 2 (полностью локально).} Без мощной GPU локальная модель 7--8B отвечает слишком медленно и нестабильно разбирает несколько действий в одной казахской фразе. Это противоречит критериям «скорость» и «качество интеграции AI». Голос при этом обрабатывается только на нашем сервере, а локальный режим LLM сохранён как опция.

\subsection{Защита от выдуманных данных}

{\tblstart
\begin{xltabular}{\textwidth}{P{4cm} L}
\toprule
\textbf{Риск} & \textbf{Механизм} \\
\midrule
\endhead
Выдуманный исполнитель & LLM возвращает только упоминание (\code{assignee\_mention}). Backend ищет его среди участников проекта: нормализация казахских падежей, транслитерация, нечёткий поиск. Не найден или найдено несколько — поле помечается, пользователь выбирает из списка \\
Выдуманный срок & LLM возвращает дату и исходную фразу (\code{deadline\_phrase}). Если фразы нет в транскрипте, дата отбрасывается. Backend проверяет дату: не в прошлом, с учётом часового пояса пользователя \\
Неверная задача при изменении & LLM возвращает описание задачи. Backend ищет кандидатов (pg\_trgm по названию в доступных проектах): 0 — «не найдено», больше 1 — уточнение у пользователя \\
Выдуманные задачи в ответах (Search, Summary) & Запрос выполняет backend по реальным данным. LLM формирует только фильтр (статус, исполнитель, срок), в ответ подставляются реальные строки из БД \\
Действие без подтверждения & Изменения выполняются только через \code{POST /ai-commands/\{id\}/confirm} \\
\bottomrule
\end{xltabular}}

% =====================================================
\section{Frontend}

{\tblstart
\begin{xltabular}{\textwidth}{P{3.6cm} L L L}
\toprule
\textbf{Критерий} & \textbf{React + TypeScript (Vite)} & \textbf{Vue 3} & \textbf{Angular} \\
\midrule
\endhead
Drag-and-drop для Kanban & dnd-kit, много готовых компонентов & vuedraggable & CDK Drag\&Drop \\
Серверное состояние и real-time & TanStack Query + WebSocket & Pinia + composables & RxJS \\
Типизация с backend & Генерация типов из OpenAPI & Аналогично & Аналогично \\
Порог входа и рынок & Самый распространённый & Простой & Высокий \\
\bottomrule
\end{xltabular}}

\textbf{Решение: React + TypeScript + Vite.} Tailwind CSS и shadcn/ui для компонентов, dnd-kit для перетаскивания карточек, TanStack Query для данных, клиент API генерируется из OpenAPI-спецификации FastAPI (единые типы на фронте и бэке), react-i18next для kk/ru/en. Голос записывается через браузерный MediaRecorder API, отдельное мобильное приложение не нужно.

% =====================================================
\section{Backend}

{\tblstart
\begin{xltabular}{\textwidth}{P{3.6cm} L L L}
\toprule
\textbf{Критерий} & \textbf{FastAPI} & \textbf{Django + DRF} & \textbf{Flask} \\
\midrule
\endhead
Асинхронность (AI-вызовы, WebSocket) & Нативная & Частичная & Через расширения \\
Валидация & Pydantic: одни и те же схемы для AI-действий, API и проверки & Serializers & Вручную \\
OpenAPI-спецификация & Автоматически & Через дополнения & Через дополнения \\
WebSocket & Встроено & Django Channels & Расширения \\
\bottomrule
\end{xltabular}}

\textbf{Решение: Python 3.12 + FastAPI}, SQLAlchemy 2.0 (async) и Alembic. Модульная структура по доменам: \code{auth}, \code{workspaces}, \code{projects}, \code{tasks}, \code{comments}, \code{notifications}, \code{ai\_commands}, \code{audit}.

\textbf{AI-воркер} — отдельный процесс (ARQ или Celery + Redis). Модели STT загружаются в память один раз, API не блокируется, воркеры масштабируются отдельно.

% =====================================================
\section{База данных и хранилища}

{\tblstart
\begin{xltabular}{\textwidth}{P{4.4cm} L L}
\toprule
\textbf{Критерий} & \textbf{PostgreSQL} & \textbf{MongoDB} \\
\midrule
\endhead
Модель данных (workspace $\rightarrow$ проект $\rightarrow$ доска $\rightarrow$ колонка $\rightarrow$ задача, участники, роли) & Реляционная, внешние ключи, транзакции & Связи вручную \\
Права доступа и целостность & Ограничения на уровне БД & В коде \\
Поиск задач по описанию из голоса & pg\_trgm; при необходимости pgvector & Atlas Search \\
Гибкие данные AI (транскрипт, план) & JSONB & Нативно \\
Аудит и история & Append-only таблица & Коллекция \\
\bottomrule
\end{xltabular}}

\textbf{Решение:}
\begin{itemize}
  \item \textbf{PostgreSQL 16} — основная БД (pg\_trgm, опционально pgvector).
  \item \textbf{Redis} — очередь AI-задач, pub/sub для WebSocket между инстансами API, rate limiting, ключи идемпотентности.
  \item \textbf{MinIO} (S3-совместимое хранилище) — голосовые сообщения; в облаке заменяется на S3 без изменения кода.
\end{itemize}

Ключевые сущности (подробно — в Database Schema): \code{users}, \code{workspaces}, \code{workspace\_members}, \code{projects}, \code{project\_members}, \code{boards}, \code{columns}, \code{tasks}, \code{task\_comments}, \code{task\_history}, \code{notifications}, \code{voice\_messages}, \code{ai\_commands}, \code{audit\_log}.

% =====================================================
\section{Безопасность и защита данных}

{\tblstart
\begin{xltabular}{\textwidth}{P{4cm} L}
\toprule
\textbf{Требование} & \textbf{Решение} \\
\midrule
\endhead
Аутентификация & JWT: короткий access-токен и refresh в httpOnly cookie; пароли хешируются Argon2 \\
Права доступа & RBAC: роли workspace (owner, admin, member) и проекта (manager, member, viewer); проверка в каждом сервисе, включая AI-действия \\
Голосовые сообщения & Приватный bucket, шифрование at rest, доступ по короткоживущим presigned URL, автоудаление через 30 дней (настраивается); транскрипт сохраняется в истории \\
Данные в облачной LLM & Только текст команды и имена участников проекта; без голоса, email и содержимого других задач. Для on-premise — локальная LLM \\
Повторное создание задач & Статус AI-команды \code{pending} $\rightarrow$ \code{confirmed} $\rightarrow$ \code{executed}; подтверждение идемпотентно, повторный клик ничего не дублирует; предупреждение о похожей недавно созданной задаче \\
Недоступность AI & Таймауты и повторы, переключение на fallback-провайдера, понятное сообщение; ручное создание задачи всегда доступно \\
История действий & \code{task\_history} (что, кем, когда, через UI или AI) и \code{audit\_log} \\
Транспорт & HTTPS/WSS, CORS, rate limiting на AI-эндпоинты \\
\bottomrule
\end{xltabular}}

% =====================================================
\section{Уведомления}

\begin{itemize}
  \item \textbf{In-app в реальном времени:} WebSocket — назначение, изменение статуса, комментарий, упоминание.
  \item \textbf{Email:} SMTP (при разработке — Mailpit) о назначении и приближающихся сроках.
  \item \textbf{Напоминания о сроках:} периодическая задача (ARQ cron или Celery beat) раз в 15 минут находит задачи со сроком в ближайшие 24 часа и просроченные.
  \item Пользователь сам настраивает, какие уведомления получать.
\end{itemize}

% =====================================================
\section{Итоговый стек}

{\tblstart
\begin{xltabular}{\textwidth}{P{3.2cm} P{6.6cm} L}
\toprule
\textbf{Слой} & \textbf{Технология} & \textbf{Ключевая причина} \\
\midrule
\endhead
Языки & kk, ru, en, смешанная речь & Требование продукта \\
Frontend & React + TypeScript + Vite, Tailwind + shadcn/ui, dnd-kit, TanStack Query, react-i18next & Зрелая экосистема для Kanban, типы из OpenAPI \\
Запись голоса & MediaRecorder API (браузер) & Без отдельного приложения \\
Backend & Python 3.12 + FastAPI, SQLAlchemy 2.0, Alembic & Async, Pydantic, автоматический OpenAPI \\
AI-воркер & ARQ / Celery + Redis & Не блокирует API, масштабируется отдельно \\
Speech-to-Text & Self-hosted: Whisper large-v3-turbo (ru/en, определение языка) + Whisper, дообученный на KSC2 (kk); облачный fallback опционально & Голос не покидает сервер; качество казахского \\
LLM & Облачная LLM Flash/mini-класса с tool calling; локальная Qwen 3 / KazLLM (Ollama) — on-premise режим и fallback & Скорость и надёжность; приватный режим сохранён \\
БД & PostgreSQL 16 + pg\_trgm (+ pgvector) & Реляционная модель, поиск задач, JSONB \\
Кэш, очереди, pub/sub & Redis & Очереди AI, WebSocket между инстансами, идемпотентность \\
Голосовые & MinIO (S3), шифрование, TTL 30 дней & Приватность, переносимость в облако \\
API & REST (OpenAPI) + WebSocket & CRUD и real-time доска \\
Auth & JWT + refresh, Argon2, RBAC & Права на workspace и проект \\
Уведомления & WebSocket + email, периодические напоминания & Изменения и сроки \\
Инфраструктура & Docker Compose (frontend, api, worker, postgres, redis, minio, ollama) & Запуск одной командой \\
Тестирование & pytest, httpx, Playwright (e2e), eval-набор AI-сценариев (30+), jiwer (WER) & Проверка и AI, и операций системы \\
CI и процесс & GitLab CI (линтеры, тесты, сборка), YouTrack & Цифровой след \\
\bottomrule
\end{xltabular}}

% =====================================================
\section{Проверка выбора}

{\tblstart
\begin{xltabular}{\textwidth}{P{3.6cm} L P{4.6cm}}
\toprule
\textbf{Что проверяем} & \textbf{Метрика} & \textbf{Набор} \\
\midrule
\endhead
STT & WER / CER по языкам отдельно для режимов auto и hint; время & 20+ записей kk/ru/en/mixed, \code{stt\_benchmark.py} \\
Определение языка & Точность определения kk/ru/en (режим auto); маршрутизация смешанных записей & Те же записи \\
Распознавание намерения & Accuracy по 8 действиям и классу «не команда» & 30+ сценариев из Test Cases \\
Извлечение полей & Точность по полям (title, assignee, deadline, priority) & Те же сценарии \\
Выдуманные данные & Доля значений, которых нет в транскрипте & Те же сценарии \\
Скорость & Время от конца записи до превью (цель — менее 5 с) & Те же сценарии \\
Сравнение LLM & Облачная и локальная на одном наборе & Подтверждение раздела~\ref{sec:llm} \\
\bottomrule
\end{xltabular}}

Результаты сохраняются в \code{benchmarks/results/} и сводятся в этот документ.

% =====================================================
\section{Риски}

{\tblstart
\begin{xltabular}{\textwidth}{P{5.4cm} L}
\toprule
\textbf{Риск} & \textbf{Митигация} \\
\midrule
\endhead
Низкая точность STT на смешанной речи & Ручной выбор языка, редактирование транскрипта до обработки \\
Недоступность облачной LLM & Fallback на локальную модель; ручное создание задач всегда работает \\
Ошибочное сопоставление исполнителя или задачи & Детерминированный поиск и подтверждение; уточнение при нескольких совпадениях \\
Нет GPU для self-hosted STT & Модели меньшего размера и int8 на CPU; облачный STT-fallback \\
Утечка голосовых данных & Приватное хранилище, шифрование, TTL, доступ по presigned URL \\
\bottomrule
\end{xltabular}}

% =====================================================
\section{Источники}

\begin{enumerate}[label={[\arabic*]}, leftmargin=*]
  \item Исследование STT/TTS для казахского языка (Whisper, GPT-4 Transcribe, ElevenLabs, Soyle). \url{https://api.crossref.org/works/10.3390%2FINFO16100879}
  \item Whisper, дообученный на KSC2: \url{https://huggingface.co/akuzdeuov/whisper-base.kk}, \url{https://huggingface.co/abilmansplus/whisper-turbo-ksc2}
  \item Whisper large-v3 для казахского (tulpar). \url{https://huggingface.co/olzhasAl/whisper-large-v3-tulpar}
  \item faster-whisper large-v3-turbo. \url{https://huggingface.co/deepdml/faster-whisper-large-v3-turbo-ct2}
  \item ISSAI Soyle. \url{https://github.com/IS2AI/Soyle}
  \item ISSAI KAZ-LLM. \url{https://issai.nu.edu.kz/kazllm/}
  \item ElevenLabs Scribe, казахский. \url{https://elevenlabs.io/speech-to-text/kazakh}
  \item ElevenLabs API pricing. \url{https://developer.puter.com/tutorials/elevenlabs-api-pricing/}
  \item OpenAI API pricing. \url{https://developers.openai.com/api/docs/pricing}
  \item Gemini Flash: structured output, tool calling, цены. \url{https://docs.apiyi.com/en/api-capabilities/gemini-3-6-flash/overview.md}
\end{enumerate}

\end{document}
